% ============================================================
\chapter{Exposants de Lyapunov}
\label{ch:exposants}
% ============================================================

\section{Introduction}

Comment mesurer le chaos~? La sensibilit\'e aux conditions initiales est une notion qualitative~: deux orbites voisines divergent, mais \`a quelle vitesse~? Alexandre Lyapunov, dans sa th\`ese de 1892, avait d\'ej\`a introduit les exposants caract\'eristiques pour \'etudier la stabilit\'e des \'equations diff\'erentielles. Mais c'est dans les ann\'ees 1960--1970, avec les travaux de Valery Oseledets (th\'eor\`eme multiplicatif ergodique, 1968), que ces exposants trouvent leur formulation d\'efinitive pour les syst\`emes dynamiques g\'en\'eraux. L'exposant de Lyapunov maximal quantifie le taux de divergence exponentielle des orbites voisines~: s'il est positif, le syst\`eme est chaotique. Le spectre complet des exposants r\'ev\`ele la g\'eom\'etrie de l'attracteur, distinguant les directions d'expansion, de contraction et de neutralit\'e.

\begin{intuition}
Imaginons deux particules tr\`es proches dans un flot. Si elles s'\'eloignent
exponentiellement vite, l'exposant de Lyapunov est positif et le syst\`eme est
chaotique. Si elles se rapprochent, l'exposant est n\'egatif et le syst\`eme
est stable. L'exposant mesure pr\'ecis\'ement ce taux exponentiel.
\end{intuition}

\section{D\'efinition pour les syst\`emes discrets}

\begin{definition}[Exposant de Lyapunov pour une application unidimensionnelle]
\label{def:lyap-1d}
Soit $f : \R \to \R$ de classe $\mathcal{C}^1$ et $x_0$ un point initial.
L'\textbf{exposant de Lyapunov}\index{exposant de Lyapunov!unidimensionnel}
en $x_0$ est
\[
    \lambda(x_0) = \lim_{n \to \infty} \frac{1}{n} \sum_{k=0}^{n-1}
    \ln \abs{f'(x_k)},
\]
o\`u $x_k = f^k(x_0)$, lorsque cette limite existe.
\end{definition}

\begin{remarque}
Par la r\`egle de la cha\^ine, $\abs{(f^n)'(x_0)} = \prod_{k=0}^{n-1} \abs{f'(x_k)}$,
donc $\lambda(x_0) = \lim_{n \to \infty} \frac{1}{n} \ln \abs{(f^n)'(x_0)}$.
Le nombre $e^\lambda$ repr\'esente le facteur moyen d'\'etirement par it\'eration.
\end{remarque}

\begin{exemple}[Application logistique]
Pour $f_r(x) = rx(1-x)$ avec $r = 4$, on sait que la mesure invariante est
$\mu(\mathrm{d}x) = \frac{1}{\pi\sqrt{x(1-x)}}\mathrm{d}x$. L'exposant de
Lyapunov vaut
\[
    \lambda = \int_0^1 \ln\abs{f_4'(x)}\,\mu(\mathrm{d}x)
    = \int_0^1 \ln\abs{4 - 8x}\,\frac{\mathrm{d}x}{\pi\sqrt{x(1-x)}} = \ln 2.
\]
\end{exemple}

\section{Exposants de Lyapunov en dimension sup\'erieure}

\begin{definition}[Matrice de Lyapunov]
Soit $\dot{x} = f(x)$ un syst\`eme en dimension $n$ et $\Phi(t, x_0)$ la
matrice fondamentale de l'\'equation variationnelle $\dot{Y} = Df(\varphi(t, x_0))Y$.
La \textbf{matrice de Lyapunov}\index{matrice de Lyapunov} est
\[
    \Lambda(x_0) = \lim_{t \to \infty} \left[\Phi(t, x_0)^\top \Phi(t, x_0)\right]^{1/(2t)}.
\]
\end{definition}

\begin{definition}[Spectre de Lyapunov]
Les \textbf{exposants de Lyapunov}\index{exposant de Lyapunov!spectre} du
syst\`eme en $x_0$ sont les logarithmes des valeurs propres de $\Lambda(x_0)$,
not\'es
\[
    \lambda_1 \geq \lambda_2 \geq \cdots \geq \lambda_n.
\]
Ils mesurent les taux de dilatation dans les $n$ directions principales.
\end{definition}

\begin{formules}[title=\textbf{Interpr\'etation g\'eom\'etrique des exposants}]
\begin{itemize}
    \item $\lambda_i > 0$ : dilatation exponentielle dans la $i$-\`eme direction.
    \item $\lambda_i = 0$ : ni contraction ni expansion (direction du flot).
    \item $\lambda_i < 0$ : contraction exponentielle.
    \item $\lambda_1 > 0$ : le syst\`eme est \textbf{chaotique}.
    \item $\sum_{i=1}^n \lambda_i < 0$ : le syst\`eme est \textbf{dissipatif}.
\end{itemize}
\end{formules}

\section{Le th\'eor\`eme ergodique multiplicatif d'Oseledets}

\begin{theoreme}[Oseledets, 1968]
\label{thm:oseledets}
\index{th\'eor\`eme d'Oseledets}
Soit $f : M \to M$ un diff\'eomorphisme pr\'eservant une mesure de probabilit\'e
ergodique $\mu$, et soit $A(x) = Df(x)$. Supposons
$\int \ln^+ \norm{A(x)} \, \mathrm{d}\mu(x) < \infty$. Alors, pour
$\mu$-presque tout $x$, il existe des r\'eels
$\lambda_1 > \lambda_2 > \cdots > \lambda_s$ et une d\'ecomposition
\[
    \R^n = E_1(x) \oplus E_2(x) \oplus \cdots \oplus E_s(x)
\]
tels que pour tout $v \in E_i(x) \setminus \{0\}$ :
\[
    \lim_{n \to \pm\infty} \frac{1}{n} \ln \norm{Df^n(x) \cdot v} = \lambda_i.
\]
De plus, $Df(x) \cdot E_i(x) = E_i(f(x))$ et $\dim E_i(x)$ est constant
$\mu$-presque partout.
\end{theoreme}

\begin{remarque}
Le th\'eor\`eme d'Oseledets g\'en\'eralise le concept d'exposants de Lyapunov
\`a tout cocycle lin\'eaire au-dessus d'un syst\`eme ergodique. Les entiers
$m_i = \dim E_i(x)$ sont les \textbf{multiplicit\'es} des exposants.
\end{remarque}

\begin{definition}[Sous-espaces d'Oseledets]
Les espaces $E_i(x)$ du th\'eor\`eme~\ref{thm:oseledets} sont appel\'es
\textbf{sous-espaces d'Oseledets}\index{sous-espaces d'Oseledets}. La
d\'ecomposition $\bigoplus E_i(x)$ est la \textbf{d\'ecomposition d'Oseledets}.
\end{definition}

\section{Propri\'et\'es fondamentales}

\begin{proposition}[Exposant le long du flot]
Pour un syst\`eme autonome continu $\dot{x} = f(x)$, la direction du flot
$f(x)$ est toujours associ\'ee \`a un exposant de Lyapunov nul. Pour un
attracteur chaotique en dimension $n$, le spectre satisfait
$\lambda_1 > 0 = \lambda_2 > \lambda_3 \geq \cdots \geq \lambda_n$.
\end{proposition}

\begin{proposition}[Formule de Liouville]
\label{prop:liouville}
La somme des exposants de Lyapunov est li\'ee \`a la divergence du champ :
\[
    \sum_{i=1}^{n} \lambda_i = \lim_{T \to \infty} \frac{1}{T}
    \int_0^T \tr\bigl(Df(\varphi(t, x_0))\bigr) \, \mathrm{d}t.
\]
Pour le syst\`eme de Lorenz ($\sigma = 10$, $\rho = 28$, $\beta = 8/3$) :
$\lambda_1 + \lambda_2 + \lambda_3 = -(\sigma + 1 + \beta) \approx -13{,}67$.
\end{proposition}

\begin{theoreme}[Formule de Pesin]
\index{formule de Pesin}
Soit $f$ un diff\'eomorphisme $\mathcal{C}^2$ d'une vari\'et\'e compacte
pr\'eservant une mesure ergodique $\mu$ absolument continue par rapport \`a
Lebesgue. Alors l'entropie m\'etrique de $f$ v\'erifie
\[
    h_\mu(f) = \sum_{\lambda_i > 0} m_i \lambda_i,
\]
o\`u la somme porte sur les exposants positifs compt\'es avec multiplicit\'e.
\end{theoreme}

\begin{attention}[title=\textbf{Exposants et dimension}]
La formule de Kaplan-Yorke relie les exposants de Lyapunov \`a la dimension
de l'attracteur : $D_{KY} = j + \frac{\sum_{i=1}^{j} \lambda_i}{\abs{\lambda_{j+1}}}$,
o\`u $j$ est le plus grand entier tel que $\sum_{i=1}^{j} \lambda_i \geq 0$.
Cette formule est conjecturalement \'egale \`a la dimension d'information.
\end{attention}

\section{Calcul num\'erique des exposants de Lyapunov}

\begin{definition}[M\'ethode de Benettin et al.]
\index{algorithme de Benettin}
L'algorithme standard pour calculer les $p$ premiers exposants de Lyapunov :
\begin{enumerate}
    \item Initialiser une base orthonorm\'ee $\{v_1, \ldots, v_p\}$.
    \item Int\'egrer simultan\'ement l'orbite et les $p$ vecteurs tangents
    sur un intervalle $\tau$.
    \item Appliquer Gram-Schmidt aux vecteurs tangents pour les
    r\'eorthogonaliser. Enregistrer les normes avant orthogonalisation.
    \item R\'ep\'eter les \'etapes 2--3 pendant $N$ it\'erations.
    \item Estimer $\lambda_i \approx \frac{1}{N\tau}\sum_{k=1}^{N} \ln \norm{v_i^{(k)}}$.
\end{enumerate}
\end{definition}

\begin{codeblock}[title=\textbf{Calcul des exposants de Lyapunov du syst\`eme de Lorenz}]
\begin{minted}{python}
import numpy as np
from scipy.integrate import solve_ivp

def lorenz(t, state, sigma=10, rho=28, beta=8/3):
    x, y, z = state[:3]
    dxdt = [sigma*(y - x), rho*x - y - x*z, x*y - beta*z]
    # Jacobian times perturbation vectors
    J = np.array([[-sigma, sigma, 0],
                  [rho - z, -1, -x],
                  [y, x, -beta]])
    Y = state[3:].reshape(3, 3)
    dYdt = J @ Y
    return dxdt + dYdt.flatten().tolist()

n = 3
x0 = [1.0, 1.0, 1.0] + list(np.eye(n).flatten())
dt, N_steps = 0.01, 10000
lyap = np.zeros(n)

for i in range(N_steps):
    sol = solve_ivp(lorenz, [0, dt], x0, max_step=dt/2)
    x0[:3] = sol.y[:3, -1]
    Y = sol.y[3:, -1].reshape(n, n)
    Q, R = np.linalg.qr(Y)
    lyap += np.log(np.abs(np.diag(R)))
    x0[3:] = Q.flatten()

lyap /= (N_steps * dt)
print(f"Exposants de Lyapunov: {lyap}")
# Resultat typique: [0.91, 0.00, -14.57]
\end{minted}
\end{codeblock}

\section{Applications \`a la d\'etection du chaos}

\begin{proposition}[Crit\`ere de chaos]
Un syst\`eme d\'eterministe est chaotique si et seulement si son plus grand
exposant de Lyapunov est strictement positif : $\lambda_1 > 0$.
\end{proposition}

\begin{exemple}[Classification par les exposants de Lyapunov]
Pour un syst\`eme tridimensionnel :
\begin{itemize}
    \item $(-, -, -)$ : point fixe stable.
    \item $(0, -, -)$ : cycle limite stable.
    \item $(0, 0, -)$ : tore quasi-p\'eriodique $T^2$.
    \item $(+, 0, -)$ : attracteur \'etrange (chaos).
\end{itemize}
\end{exemple}

\begin{lemme}[Temps de pr\'edictibilit\'e]
Le temps de Lyapunov $\tau_L = 1/\lambda_1$ d\'efinit l'\'echelle de temps
caract\'eristique au-del\`a de laquelle la pr\'ediction devient impossible.
Apr\`es un temps $t$, l'incertitude initiale $\delta_0$ cro\^it selon
$\delta(t) \sim \delta_0 \, e^{\lambda_1 t}$.
\end{lemme}

\begin{exemple}[Pr\'evision m\'et\'eorologique]
Pour l'atmosph\`ere, $\lambda_1 \approx 1/(2\text{ jours})$. Une erreur initiale
de $10^{-6}$ atteint l'\'echelle synoptique ($\sim 1$) apr\`es environ
$t = \frac{\ln(10^6)}{\lambda_1} \approx 28$ jours, ce qui fixe la limite
th\'eorique de la pr\'evision m\'et\'eorologique.
\end{exemple}

\section{Exposants conditionnels et applications}

\begin{definition}[Exposants de Lyapunov conditionnels]
\index{exposants conditionnels}
Pour deux syst\`emes coupl\'es $\dot{x} = f(x)$ et $\dot{y} = g(x, y)$,
les exposants de Lyapunov \textbf{conditionnels} du syst\`eme esclave $y$
sont d\'efinis par la lin\'earisation
$\dot{\delta y} = D_y g(x(t), y(t)) \cdot \delta y$.
Ils d\'eterminent si la synchronisation $y(t) \to x(t)$ est stable.
\end{definition}

\begin{theoreme}[Synchronisation chaotique]
Deux syst\`emes chaotiques identiques coupl\'es se synchronisent si et
seulement si tous les exposants de Lyapunov conditionnels transverses sont
strictement n\'egatifs.
\end{theoreme}

\section{Exercices}

\begin{exercice}[Exposant de la tente]
Calculer analytiquement l'exposant de Lyapunov de l'application de la tente
$T(x) = 1 - \abs{2x - 1}$ sur $[0,1]$ par rapport \`a la mesure de Lebesgue.
\end{exercice}

\begin{exercice}[Spectre du chat d'Arnold]
Soit l'application du chat d'Arnold $\begin{pmatrix} x_{n+1} \\ y_{n+1} \end{pmatrix}
= \begin{pmatrix} 2 & 1 \\ 1 & 1 \end{pmatrix} \begin{pmatrix} x_n \\ y_n \end{pmatrix}
\pmod{1}$. Calculer les deux exposants de Lyapunov et v\'erifier que leur somme
est nulle (syst\`eme conservatif).
\end{exercice}

\begin{exercice}[Dimension de Kaplan-Yorke]
Pour le syst\`eme de Lorenz avec les param\`etres standard, les exposants
sont approximativement $\lambda_1 \approx 0{,}91$, $\lambda_2 = 0$,
$\lambda_3 \approx -14{,}57$. Calculer la dimension de Kaplan-Yorke.
\end{exercice}

\begin{exercice}[Convergence num\'erique]
Impl\'ementer l'algorithme de Benettin pour l'application de H\'enon
$H(x,y) = (1 - 1{,}4x^2 + y, \; 0{,}3x)$ et \'etudier la convergence
des exposants en fonction du nombre d'it\'erations.
\end{exercice}

\begin{exercice}[Formule de Pesin]
V\'erifier num\'eriquement la formule de Pesin pour l'application du chat
d'Arnold en calculant ind\'ependamment l'entropie m\'etrique et la somme des
exposants positifs.
\end{exercice}
